% Einheit 4 von 5 – Sachaufgaben (bank/lineare-gleichungssysteme/e4.jsonl, zusammenbau v0.1)
%% TODO Zweigzeile Teil 1: was hier gelernt wird (Satz in Schülersprache aus dem Einheitstitel) – Einheit 4
\einheitenkopf[e4]{Einheit 4 von 5 · Sachaufgaben}
\zweigzeile{ab Kl. 9 (am Gymnasium ab Kl. 8) · P10}
\setcounter{aufgabe}{22}

%% TODO \verfahren{…}: Verfahrensüberschrift in Sachsprache fehlt in der Bank (2.3 b)
%% TODO Titel: Ich-kann-Satz fehlt in der Bank (Platzhalter: Kettenname) – Nr. 23
%% TODO Anweisung: ein Satz über der ersten Teilaufgabe fehlt in der Bank – Nr. 23
\begin{aufgabe}{Sachaufgaben}
\begin{teile}
\teil Bäckerei: $3$ Brötchen und $2$ Brezeln kosten $3{,}10\,€$; $1$ Brötchen und $4$ Brezeln kosten $3{,}70\,€$. Was ist unbekannt? Benenne die gesuchten Größen mit $x$ und $y$ und schreib heraus, welche Zahlen rechts vom Gleichheitszeichen und welche vor $x$ und $y$ stehen.
\end{teile}
%% TODO Darstellung für schwach fehlt in der Bank (lineare-gleichungssysteme-e4-k1-s1-v1; Abschnitt „Für schwache Schüler“)
\swz[3]{Ein Eis und ein Kakao kosten zusammen $4{,}60\,€$. Drei Eis und vier Kakao kosten $15{,}80\,€$. $x$: Preis eines Eises, $y$: Preis eines Kakaos (in €). Gleichungen? Stelle I und II auf.}{}
%% TODO Darstellung für schwach fehlt in der Bank (lineare-gleichungssysteme-e4-k1-s1-v2; Abschnitt „Für schwache Schüler“)
\swz[3]{Ein Buch und ein Kalender kosten zusammen $21\,€$. Zwei Bücher und fünf Kalender kosten $69\,€$. $x$: Preis eines Buchs, $y$: Preis eines Kalenders (in €). Gleichungen? Stelle I und II auf.}{}
%% TODO Darstellung für schwach fehlt in der Bank (lineare-gleichungssysteme-e4-k1-s1-v3; Abschnitt „Für schwache Schüler“)
\swz[3]{Ein Liter Milch und ein Brot kosten zusammen $3{,}40\,€$. Vier Liter Milch und zwei Brote kosten $8{,}80\,€$. $x$: Preis eines Liters Milch, $y$: Preis eines Brots (in €). Gleichungen? Stelle I und II auf.}{}
%% TODO Darstellung für schwach fehlt in der Bank (lineare-gleichungssysteme-e4-k1-s1-v4; Abschnitt „Für schwache Schüler“)
\swz[3]{Ein Fußball und ein Basketball kosten zusammen $45\,€$. Zwei Fußbälle und drei Basketbälle kosten $110\,€$. $x$: Preis eines Fußballs, $y$: Preis eines Basketballs (in €). Gleichungen? Stelle I und II auf.}{}
%% TODO Darstellung für schwach fehlt in der Bank (lineare-gleichungssysteme-e4-k1-s1-v5; Abschnitt „Für schwache Schüler“)
\swz[3]{Eine Pizza und ein Salat kosten zusammen $14\,€$. Zwei Pizzen und fünf Salate kosten $46\,€$. $x$: Preis einer Pizza, $y$: Preis eines Salats (in €). Gleichungen? Stelle I und II auf.}{}
\swfrage{Was bleibt gleich, was ändert sich?}
%% TODO Darstellung für schwach fehlt in der Bank (lineare-gleichungssysteme-e4-k1-s2-v1; Abschnitt „Für schwache Schüler“)
\swz[3]{Freibad: Am Morgen wurden $50$ Karten verkauft, Kinderkarten zu $2{,}80\,€$ und Erwachsenenkarten zu $4{,}50\,€$, zusammen für $177{,}40\,€$. $x$: Anzahl der Kinderkarten, $y$: Anzahl der Erwachsenenkarten. Gleichungen? Stelle I und II auf.}{}
%% TODO Darstellung für schwach fehlt in der Bank (lineare-gleichungssysteme-e4-k1-s3-v1; Abschnitt „Für schwache Schüler“)
\swz[3]{Jugendherberge: $23$ Zimmer, Zweibett- und Vierbettzimmer, mit zusammen $68$ Betten. $x$: Anzahl der Zweibettzimmer, $y$: Anzahl der Vierbettzimmer. Gleichungen? Stelle I und II auf.}{}
%% TODO Darstellung für schwach fehlt in der Bank (lineare-gleichungssysteme-e4-k1-s4-v1; Abschnitt „Für schwache Schüler“)
\swz[3]{Die Summe zweier Zahlen ist $31$, ihre Differenz ist $7$. $x$: die größere Zahl, $y$: die kleinere Zahl. Gleichungen? Stelle I und II auf.}{}
%% TODO Darstellung für schwach fehlt in der Bank (lineare-gleichungssysteme-e4-k1-s5-v1; Abschnitt „Für schwache Schüler“)
\swz[3]{Schreibwaren: Ein Block und ein Füller kosten zusammen $8\,€$, vier Blöcke und zwei Füller $26\,€$. Dazu gehört I: $x + y = 8$, II: $4x + 2y = 26$. Wofür stehen $x$ und $y$?}{}
%% TODO Darstellung für schwach fehlt in der Bank (lineare-gleichungssysteme-e4-k1-s6-v1; Abschnitt „Für schwache Schüler“)
\swz[3]{Obstkorb: $x$ ist die Anzahl der Äpfel, $y$ die Anzahl der Birnen im Korb. Was sagt $x + y = 25$? Formuliere einen Satz.}{}
%% TODO Darstellung für schwach fehlt in der Bank (lineare-gleichungssysteme-e4-k1-s7-v1; Abschnitt „Für schwache Schüler“)
\swz{$\text{Imbiss: Zwei Burger und eine Portion Pommes kosten $14{,}50\,€$, ein Burger und zwei Portionen Pommes $12{,}50\,€$. Was kostet ein Burger, was eine Portion Pommes? Stelle ein System auf, löse es und antworte.}$}{}
\swz[3]{Café: Ein Croissant und ein Muffin kosten zusammen $4{,}10\,€$. Fünf Croissants und zwei Muffins kosten $13{,}60\,€$. $x$ ist der Preis eines Croissants, $y$ der Preis eines Muffins in Euro. Stelle ein Gleichungssystem auf. \hfill (P10 2024 OS)}{}
\end{aufgabe}

%% TODO \verfahren{…}: Verfahrensüberschrift in Sachsprache fehlt in der Bank (2.3 b)
%% TODO Titel: Ich-kann-Satz fehlt in der Bank (Platzhalter: Kettenname) – Nr. 24
%% TODO Anweisung: ein Satz über der ersten Teilaufgabe fehlt in der Bank – Nr. 24
\begin{aufgabe}{Sachaufgaben, Sek II}
\begin{teile}
\teil Was ist unbekannt? „Ein Rechteck hat den Umfang $30$ cm; die eine Seite $x$ ist $3$ cm länger als die andere.“ Kreuze an, wofür $x$ steht.\\ \kreuz{eine Länge}\\ \kreuz{eine Anzahl}\\ \kreuz{ein Preis}\\ \kreuz{ein Anteil}
\end{teile}
%% TODO Darstellung für schwach fehlt in der Bank (lineare-gleichungssysteme-e4-k2-s1-v1; Abschnitt „Für schwache Schüler“)
\swz[3]{Bäckerei: $x$ ist die Zahl der Roggenbrote, $y$ die Zahl der Weizenbrote. I: $x + y = 60$, II: $x = 2y$. Was sagen I und II? Formuliere je einen Satz.}{}
%% TODO Darstellung für schwach fehlt in der Bank (lineare-gleichungssysteme-e4-k2-s1-v2; Abschnitt „Für schwache Schüler“)
\swz[3]{Kino: Verkauft wurden $x$ Karten zu $8\,€$ und $y$ Karten zu $14\,€$. I: $x + y = 150$, II: $8x + 14y = 1\,500$. Was sagen I und II? Formuliere je einen Satz.}{}
%% TODO Darstellung für schwach fehlt in der Bank (lineare-gleichungssysteme-e4-k2-s1-v3; Abschnitt „Für schwache Schüler“)
\swz[3]{Garten: $x$ kg Erde und $y$ kg Sand werden gemischt. I: $x + y = 50$, II: $x - y = 30$. Was sagen I und II? Formuliere je einen Satz.}{}
%% TODO Darstellung für schwach fehlt in der Bank (lineare-gleichungssysteme-e4-k2-s1-v4; Abschnitt „Für schwache Schüler“)
\swz[3]{Rechteck: Die Seiten sind $x$ cm und $y$ cm lang. I: $2x + 2y = 36$, II: $x = y + 4$. Was sagen I und II? Formuliere je einen Satz.}{}
%% TODO Darstellung für schwach fehlt in der Bank (lineare-gleichungssysteme-e4-k2-s1-v5; Abschnitt „Für schwache Schüler“)
\swz[3]{Alter: Paul ist $x$ Jahre alt, seine Mutter $y$ Jahre. I: $y = 3x$, II: $x + y = 48$. Was sagen I und II? Formuliere je einen Satz.}{}
\swfrage{Was bleibt gleich, was ändert sich?}
%% TODO Darstellung für schwach fehlt in der Bank (lineare-gleichungssysteme-e4-k2-s2-v1; Abschnitt „Für schwache Schüler“)
\swz[3]{Tierfutter: Drei Sorten enthalten $15\,\%$, $25\,\%$ und $45\,\%$ Eiweiß. Eine Mischung aus den Anteilen $u$, $v$ und $w$ der drei Sorten soll $30\,\%$ Eiweiß haben: I: $u + v + w = 1$, II: $15u + 25v + 45w = 30$. Was bedeuten I und II?}{}
\swz[3]{Aus drei Düngern mit $4\,\%$, $8\,\%$ und $16\,\%$ Stickstoff wird ein Dünger mit $9\,\%$ Stickstoff gemischt. Dazu gehört I: $4u + 8v + 16w = 9$, II: $u + v + w = 1$. Interpretiere das Gleichungssystem im Sachzusammenhang. \hfill (Abitur 2024 LK)}{}
\end{aufgabe}

%% TODO Titel: Ich-kann-Satz fehlt in der Bank (Platzhalter: Kettenname) – Nr. 25
%% TODO Anweisung: ein Satz über der ersten Teilaufgabe fehlt in der Bank – Nr. 25
\begin{aufgabe}{Mischungen und Alter (Vorrat)}
%% TODO Darstellung für schwach fehlt in der Bank (lineare-gleichungssysteme-e4-k3-s1-v1; Abschnitt „Für schwache Schüler“)
\swz{$\text{Studentenfutter: Nüsse zu $14\,€$ je kg und Rosinen zu $5\,€$ je kg sollen $30$ kg Mischung zu $8\,€$ je kg ergeben. Wie viel kg von jeder Sorte? Stelle ein System auf und löse es.}$}{}
\end{aufgabe}

%% TODO Titel: Ich-kann-Satz fehlt in der Bank (Platzhalter: Kettenname) – Nr. 26
%% TODO Anweisung: ein Satz über der ersten Teilaufgabe fehlt in der Bank – Nr. 26
\begin{aufgabe}{Probe im Text (Rechnung mit den Zahlen der Aufgabe)}
%% TODO Darstellung für schwach fehlt in der Bank (lineare-gleichungssysteme-e4-k4-s1-v1; Abschnitt „Für schwache Schüler“)
\swz[3]{Probe im Text: „Eine Tasse und ein Teller kosten zusammen $9\,€$; drei Tassen und zwei Teller kosten $23\,€$.“ Lisa findet: Tasse $5\,€$, Teller $4\,€$. Prüfe mit den Zahlen aus dem Text.}{}
\end{aufgabe}

%% TODO Titel: Ich-kann-Satz fehlt in der Bank (Platzhalter: Kettenname) – Nr. 27
%% TODO Anweisung: ein Satz über der ersten Teilaufgabe fehlt in der Bank – Nr. 27
\begin{aufgabe}{Bestandteile einer Gleichung deuten}
%% TODO Darstellung für schwach fehlt in der Bank (lineare-gleichungssysteme-e4-k5-s1-v1; Abschnitt „Für schwache Schüler“)
\swz[3]{Handyrechnung: $x$ ist die Zahl der Minuten, $y$ die Zahl der SMS. Was bedeutet in $0{,}09x + 0{,}19y = 12{,}60$ die Zahl $0{,}09$, was die Zahl $12{,}60$?}{}
\end{aufgabe}

%% TODO Titel: Ich-kann-Satz fehlt in der Bank (Platzhalter: Kettenname) – Nr. 28
%% TODO Anweisung: ein Satz über der ersten Teilaufgabe fehlt in der Bank – Nr. 28
\begin{aufgabe}{Sachaufgaben – Fehler finden, Begründen}
\swz[3]{„Ein Kakao und ein Muffin kosten zusammen $4{,}20\,€$; zwei Kakao und fünf Muffins kosten $14{,}10\,€$.“ Ole stellt auf: I: $x + y = 14{,}10$, II: $2x + 5y = 4{,}20$. Finde den Fehler.}{}
\swz[3]{„Ein Hotel hat $17$ Zimmer, Zweibett- und Dreibettzimmer, mit zusammen $43$ Betten.“ $x$ ist die Anzahl der Zweibettzimmer, $y$ die Anzahl der Dreibettzimmer. Kim rechnet $x = 8$, $y = 9$ und antwortet: „Es gibt $8$ Dreibettzimmer und $9$ Zweibettzimmer.“ Finde den Fehler.}{}
\swz[3]{Müsli: Drei Sorten mit $5\,\%$, $15\,\%$ und $30\,\%$ Nussanteil werden zu einem Müsli mit $25\,\%$ Nussanteil gemischt: I: $u + v + w = 1$, II: $5u + 15v + 30w = 25$. Nils schreibt: „$u$, $v$ und $w$ sind die Nussanteile der drei Sorten in Prozent.“ Finde den Fehler.}{}
\swfrage{Warum reicht die Angabe „Ein Brot und ein Kuchen kosten zusammen $8\,€$“ nicht, um beide Preise zu bestimmen? Begründe.}
\swfrage{Eine Mischung aus drei Sorten hat die Anteile $u$, $v$ und $w$. Warum gilt $u + v + w = 1$, und warum ist die Gleichung für den Gehalt einer Zutat eine Bilanz? Begründe.}
\end{aufgabe}

\uebersichtskasten{Zwei Unbekannte, zwei Gleichungen: x und y benennen (mit Einheit), aus jeder Angabe im Text eine Gleichung machen. \\ „Ein Heft und ein Stift kosten zusammen 2,50 €; vier Hefte und drei Stifte kosten 8,50 €.“ \quad x Preis eines Hefts, y Preis eines Stifts: I x + y = 2,50 \quad II 4x + 3y = 8,50 \quad Lösung x = 1,00 €, y = 1,50 € \\ Anzahl-und-Preis: die Anzahlen stehen vor x und y, der Gesamtbetrag steht rechts; „zusammen n Stück“ gibt x + y = n. \\ Antwort: die Lösung den Größen zuordnen und den Satz aus der Frage bilden; Probe mit den Zahlen aus dem Text. \\ Gleichung deuten: Vorzahl = Anzahl (oder Preis je Stück), Variable = das Gesuchte, rechte Seite = Gesamtwert. „x + y = 20“ heißt: zusammen zwanzig Stück.}
